% !TeX root = ../../金融经济学_熊和平_习题.tex

\xhpthoughts
\begin{problem}
\problemno{1} 什么是 Arrow-Debreu 证券市场？为什么要讨论 Arrow-Debreu 证券？
\end{problem}
\begin{problem}
\problemno{2} Arrow-Debreu 证券与现实生活中的什么金融产品类似？
\end{problem}
\begin{problem}
\problemno{3} 金融文献中在分析问题时经常假设市场是完全的，试说明完全市场的优点是什么？
\end{problem}
\begin{problem}
\problemno{4} Arrow-Debreu 证券市场的一般均衡分析分哪几步进行？
\end{problem}

\xhpexercises
\begin{problem}
\problemno{1} 考虑三状态经济中有三个证券，价格分别为 \num{5} 元、\num{5} 元和 \num{6} 元，期末的支付矩阵为
\[
    X=\begin{bmatrix}
        10 & 20 & 30\\
        10 & 10 & 10\\
        10 & 5 & 5
    \end{bmatrix}.
\]
\begin{enumerate}[label=(\arabic*)]
    \item 将三个证券表示为 A-D 证券的组合；
    \item 计算三个不同的状态证券价格。
\end{enumerate}
\end{problem}
\begin{problem}
\problemno{2} 给定下列支付矩阵，问 $k$ 为何值时市场是完全的：
\[
    X=\begin{bmatrix}
        3 & 2 & 3\\
        1 & k & 5\\
        2 & 1 & 7
    \end{bmatrix}.
\]
\end{problem}
\begin{problem}
\problemno{3} 市场上有三个证券，其价格和期末支付分别为
\[
    S=\begin{bmatrix}6\\4\\10\end{bmatrix},\qquad
    X=\begin{bmatrix}
        4 & 3 & 2\\
        2 & 7 & 3\\
        1 & 4 & 3
    \end{bmatrix}.
\]
\begin{enumerate}[label=(\arabic*)]
    \item 证明不存在套利机会；
    \item 确定状态价格向量；
    \item 确定消费计划 $(1,0,3)$ 的无套利价值；
    \item 确定以第一个证券为标的、执行价格为 $2$ 的看涨期权价格。
\end{enumerate}
\end{problem}
\begin{problem}
\problemno{4} 假定市场上有两种证券：证券 1 和证券 2，价格分别为 \num{13} 元和 \num{10} 元，两种证券期末的支付如下：
\begin{center}
\begin{tabular}{ccc}
    & 状态证券 1 的支付 & 状态证券 2 的支付\\
    \hline
    证券 1 & 11 & 0\\
    证券 2 & 2 & 20
\end{tabular}
\end{center}
假设经济人的效用函数为 $U(w_1,w_2)=w_1 w_2$，其中 $w_1,w_2$ 分别为期末不同状态下的财富水平。该经济人初始财富水平为 \num{420} 元。为了最大化效用，在期初构造投资组合时分别购买多少证券 1 和证券 2？
\end{problem}
\begin{problem}
\problemno{5} 假若期末只有一个状态发生，即期末是确定性的。一个消费者的两期消费的效用函数为 $u(c_0,c_1)=c_0 c_1^{1/2}$，其中 $c_0,c_1$ 分别为期初和期末的消费；该消费者的收入流为 $y_0=\num{2000}$ 元，$y_1=\num{1296}$ 元，且市场利率为 $8\%$。求该消费者的最优消费计划。
\end{problem}
\begin{problem}
\problemno{6} 继续考虑期末确定情形，经济人 1 在 0 期的禀赋为 $100$ 而在 1 期的禀赋为 $1$，其效用函数为 $u(c_0,c_1)=\ln c_0+\rho\ln c_1$；市场上只有一个无风险证券，无风险利率为 $r_f$。
\begin{enumerate}[label=(\arabic*)]
    \item 若 $\rho=0.5$，求经济人的投资需求；
    \item 若经济人 1 是唯一的参与者，请描述市场出清条件；
    \item 求均衡利率；
    \item 求均衡利率与时间偏好率 $\rho$ 的关系。
\end{enumerate}
\end{problem}
\begin{problem}
\problemno{7} 在第 4 题的基础上加入经济人 2，其效用函数与经济人 1 相同，禀赋为 $(1,100)$，即期初禀赋为 $1$，期末禀赋为 $100$。
\begin{enumerate}[label=(\arabic*)]
    \item 求经济人 2 的最优消费-投资策略；
    \item 求新的经济下的均衡利率，并比较单个经济人参与的经济下的均衡利率。
\end{enumerate}
\end{problem}
\begin{problem}
\problemno{8} 假设期末有三种等可能的状态 $a$、$b$ 和 $c$ 发生，市场是完全的，有两个经济人。经济人 1 的禀赋是 $100$ 个单位期初的消费品；经济人 2 的禀赋为期末的消费品，三个状态下分别为 $200$、$100$ 和 $50$。他们的效用函数相同，均为
\[
    u(c_0;c_a,c_b,c_c)=-\frac{1}{c_0}-\frac{1}{3}\Bigl(\frac{1}{c_a}+\frac{1}{c_b}+\frac{1}{c_c}\Bigr).
\]
\begin{enumerate}[label=(\arabic*)]
    \item 求解各自优化问题，得到每个人的最优消费-投资；
    \item 求均衡条件下三种状态的状态价格；
    \item 分析每个经济人的总消费与各自禀赋的关系。
\end{enumerate}
\end{problem}
